% page 600 du fac-simile (image p-599 du scan)
% bloc 152.4 x 254.9 mm ; marge gauche 8.0 mm
\pgc{-12.700mm}{0.000mm}{
\l{\cel{3}592.}
\l{}
\l{\sou{travestiar}. (\sou{trans}.) I. Metar (sur su, sur ulu) vesto, kos-}
\l{\cel{3}tumo, qua ne esas olta di lua klaso sociala o di lua se-}
\l{\cel{3}xuo. - II. Metar (sur ulu, sur su) kostumo qua igas des-}
\l{\cel{3}facila lua rekonoceso, - DEFI.}
\l{}
\l{\sou{tre}.\cel{2}(\sou{adverbo}). Alta-grade, pri lua qualeso. F.}
\l{}
\l{\sou{trefo}. (\sou{karto-ludo}). Un ek la emblemi nigra-kolora qua re-}
\l{\cel{3}prezentas folio di trifolio. - DFR.}
\l{}
\l{\sou{treko}. Mala voyeto quan facis nevole la personi, o trupi,}
\l{\cel{3}o veturi, qui pasis ofte sur la traci di lia pre-irinto,}
\l{\cel{3}en sulo naturala. - E.}
\l{}
\l{\sou{trelicho}. Telo grosiera, por la faco di varo-saki, labor-}
\l{\cel{3}vesti. - DFI.}
\l{}
\l{\sou{treliso}. (\sou{tekn.})\cel{2}Plako plu o min granda areale, ek fer-fili}
\l{\cel{3}interplektita tale ke li formacas mashi plu o min larja,}
\l{\cel{3}quan onu muntas avan fenestro o lokizas avan la fairuyo}
\l{\cel{3}di kameno, od uzas (larja-bendo) kom klozilo di gardeno,}
\l{\cel{3}di tereno, e c. EF.}
\l{}
\l{\sou{tremar}.\cel{2}(\sou{netrans}.) (\sou{pro,}\cel{1}de). Agiteso da serio de mikra}
\l{\cel{3}ocili, quin produktas la koldeso di atmosfero, febro,}
\l{\cel{3}pavoro, timo, e c. en la homo-korpo. -(\sou{anke metaf.})}
\l{\cel{3}- EFI.}
\l{}
\l{\sou{tremao.}\cel{1}Signo ek du punti inter-apuda, quan onu lokizas ho-}
\l{\cel{3}rizontale super vokalo por modifikar la pronunco di ol-}
\l{\cel{3}ca (en la lingui Germana o Germanida) o por indikar ke}
\l{\cel{3}du vokali inter-apuda ne pronuncesas diftonge (en la lin-}
\l{\cel{3}gui Franca, Angla, e c.)\cel{2}(Kom ex.:\cel{2}F. noël, poële -}
\l{\cel{3}A. coöperation, G. Brüder). - DFS.}
\l{}
\l{\sou{tremolo}. (\sou{muziko}).Tremo sencesa sur noto muzikala. - DEFIRS}
\l{}
\l{\sou{trempar}. (\sou{trans., en)}. Plunjar en liquido, por imbibar ye ta}
\l{\cel{3}liquido (la plunjajo). - F.}
\l{}
\l{\sou{tremulo}. (\sou{bot.)}\cel{2}Speco di poplo, di qua la folii, lejera,}
\l{\cel{3}tremas efike da vento mem febla. - L.\cel{1}\sou{populus tremula}.}
\l{}
\l{}
\l{\sou{treno}. Ensemblo de la vagoni quin tranas lokomotivo. - EFIS}
\l{}
\l{\sou{trepano}. I. (\sou{kirurgio)}\cel{2}Instrumento por borar la osti, pre-}
\l{\cel{3}cipue olti di la kranio, por posibligar la ekfluo de}
\l{\cel{3}puso o la retrodukto di osto profunde sinkita a lua pla-}
\l{\cel{3}so normala. - II. (\sou{tekn.})\cel{2}Granda mashino por perfora}
\l{\cel{3}sulo roka, por facar putei, serche di petrolo, e c. -}
\l{\cel{3}DEFIRS.}
\l{}
\l{\sou{trepidar.}\cel{1}(\sou{netrans.}) I. (\sou{aludante navo, automobilo, qua fun}-}
\l{\cel{3}\sou{cionas}). Tremar ye sukusi bruska quin produktas la vaor-}
\l{\cel{3}mashino, la motoro. - II. Manifestar, malgre su, signi}
\l{\cel{3}di sua nepacienteso. - EFIRS.}
}
